%======================================================================
\section{Our method, step by step}
%======================================================================

\begin{figure}[h]
\centering
\resizebox{\linewidth}{!}{%
\begin{tikzpicture}[font=\sffamily\small,>=Stealth,
  box/.style={draw=black!55,rounded corners=2pt,fill=#1,align=center,inner sep=4pt,minimum height=1.0cm,text width=2.4cm},
  box/.default=white, arr/.style={->,thick,black!65}]
\node[box=black!4] (in) at (0,0) {Photo $I$};
\node[box=blue!8] (d1) at (3.2,0.9) {Landmarker on $I$};
\node[box=blue!8] (d2) at (3.2,-0.9) {Landmarker on flip$(I)$};
\node[box=blue!15] (mt) at (6.5,0) {\textbf{Mirror test}\\\scriptsize $\delta$ and average $\bar L$};
\node[box=green!10] (map) at (9.8,0) {478$\to$68 map};
\node[box=green!18] (ax) at (13.1,0) {\textbf{Fitted axis}\\\scriptsize $A$, $A_r$};
\node[box=orange!18] (gate) at (9.8,-2.2) {\textbf{Quality gate}\\\scriptsize size, blur, light, yaw, pitch, $\delta$};
\node[box=black!8] (out) at (16.4,-1.1) {\textbf{Output}\\\scriptsize accept / caution / reject};
\draw[arr] (in) -- (d1); \draw[arr] (in) -- (d2);
\draw[arr] (d1) -- (mt); \draw[arr] (d2) -- (mt);
\draw[arr] (mt) -- (map); \draw[arr] (map) -- (ax);
\draw[arr] (mt) |- (gate); \draw[arr] (ax) -- (out); \draw[arr] (gate) -| (out);
\end{tikzpicture}}
\caption{The full system (variant M3). The numbered steps below follow this picture.}
\end{figure}

\subsection{Step 1: detect landmarks (\texttt{detector.py})}
\textbf{What:} run MediaPipe Face Landmarker on the photo. We get 478 points in pixel
coordinates and the transformation matrix, from which we read yaw, pitch and roll.\\
\textbf{Why MediaPipe:} it is pretrained (no training needed), fast on a CPU, and its mesh has
a \emph{fixed topology}: vertex 33 is always the same anatomical spot, on every face. That
is what makes fixed pairs possible.\\
\textbf{Limitation:} it gives no confidence per landmark (the \texttt{visibility} field is empty).
The mirror test (Step 5) fills this gap.

\subsection{Step 2: find each point's mirror twin (\texttt{mesh.py})}\label{sec:perm}
\textbf{What:} take the canonical face model, flip it ($x\to -x$), and for each vertex find the
vertex that lands in its place. Because the canonical model is exactly symmetric, this
match is perfect (error 0). It gives a \textbf{mirror permutation} $\pi$: $\pi(i)$ is the
twin of vertex $i$. Result: 220 bilateral pairs and 28 midline vertices (which map to
themselves), plus the iris points, paired by hand.\\
\textbf{Why:} the twin is fixed by anatomy, not guessed per photo. This removes the
baseline's wrong-partner problem completely (0\% wrong partners).

\subsection{Step 3: map the 478-point mesh to the 68-point layout}
\textbf{What:} for each of the 68 standard points, choose the mesh vertex that best plays its
role. We do this \emph{on training images only}: for every standard point $k$ and every mesh
vertex $v$, we compute how far $v$ typically is from the expert's point $k$ (median over
810 training faces, divided by IOD). Call this $E_{kv}$. For a bilateral pair $(a,b)$ of the
68 points we choose
\[ v^\star = \arg\min_v \; E_{a v} + E_{b\,\pi(v)}, \qquad a \to v^\star, \quad b \to \pi(v^\star). \]
\textbf{Why this way:} choosing the two sides \emph{together}, as mirror twins, guarantees
the mapping is itself perfectly symmetric. If the left eye corner were mapped slightly
differently from the right eye corner, the mapping alone would create fake asymmetry.\\
\textbf{Why 68 points at all:} expert annotations exist for these points, so we can check our
measurements against experts. The fitted map is also more accurate than the commonly used
hand-made map: NME \nmeFitInner\% vs.\ \nmeHandInner\% (inner points), and \nmeFitJaw\% vs.\
\nmeHandJaw\% on the jaw.

\subsection{Step 4: fit the symmetry axis and measure asymmetry (\texttt{geometry.py})}
\textbf{What:} instead of assuming the midline is vertical, we search for the line (angle
$\theta$, offset $\rho$) that best explains the face as a mirror image:
\[ (\theta^\star,\rho^\star)=\arg\min_{\theta,\rho}\sum_{i}\big\|\mathcal{R}_{\theta,\rho}(p_{\pi(i)})-p_i\big\|^2 . \]
In words: reflect each point's twin across a candidate line, measure how far it lands from
the point, and pick the line that makes these gaps as small as possible overall. We start
from a closed-form guess (midpoints and directions of the pairs) and refine it with the
Nelder--Mead optimiser. Only the 51 inner-face points are used to fit it.

Then, with residual $r_i$ = reflection of the twin minus the point:
\[ a_i=\frac{\|r_i\|}{2\cdot\text{IOD}}, \qquad A=\text{mean of } a_i \text{ over the 51 inner points}, \qquad A_r=\text{mean over region } r. \]
The factor 2 is there because a gap $r_i$ is shared by two points. $a_i$ is how far point
$i$ is from the closest perfectly symmetric face.

\textbf{Why fit the axis:}
\begin{itemize}
\item It is \textbf{rotation-invariant}: tilt the photo and $A$ does not change (a unit test checks this for $-30^\circ$ to $+25^\circ$, with scaling and shifting).
\item It does not depend on two special points being correct.
\item It needs no second detection pass, unlike the baseline's align-and-re-detect step.
\end{itemize}

\begin{idea}[Worked example (this is one of our unit tests)]
Take a perfectly symmetric face, so $A=0$. Move the inner corner of one eye 2 pixels outward.
The fitted axis barely moves, the gap between that corner and its reflected twin becomes
about 2 pixels, and each of the two corners gets $a_i \approx 1$ pixel / IOD. The
\emph{Eyes} region now has the largest regional value, as it should.
\end{idea}

\textbf{Why the jaw is excluded from $A$:} the 17 jaw points trace the visible outline
(silhouette) of the face. Their position along the cheek shifts with head pose, and
MediaPipe's outline does not correspond to them point for point (jaw NME \nmeFitJaw\% vs.\
\nmeFitInner\% for the inner points). The jaw is still shown as a region.

\subsection{Step 5: the mirror test}
\textbf{What:} run the detector a second time, on the horizontally flipped photo. Then map
the result back:
\begin{itemize}
\item a point at $x'$ in the flipped photo is at $W-x'$ in the original ($W$ = image width;
MediaPipe uses continuous coordinates, so it is $W-x'$, not $W-1-x'$, which we checked:
using $W-1$ left a systematic 1-pixel offset);
\item flipping swaps sides, so the vertex the detector calls $j$ in the flipped photo is the
original's vertex $\pi(j)$.
\end{itemize}
This gives a second set of landmarks $U$ for the \emph{same} face. A perfect detector would give
$U=L$.

\textbf{What we get from it:}
\begin{enumerate}
\item \textbf{Mirror inconsistency} $\delta$ = mean distance between $L$ and $U$ / IOD. Large
$\delta$ means the detector is unsure of the geometry. It correlates with true landmark
error ($\rho$ = \rhoMirrorNme{}), and it needs no labels.
\item \textbf{Averaged landmarks} $\bar L=\tfrac12(L+U)$. If the detector systematically
pushes points to one side, the flip pushes them the other way, and the average cancels it.
This reduces score drift on degraded photos by \driftReduction\%.
\item \textbf{Stability band} $|A(L)-A(U)|$, shown to the user as ``$\pm$''.
\item \textbf{Per-region inconsistency} $\delta_r$, used to grey out unreliable regions.
\end{enumerate}

\begin{say}
``Flipping a photo cannot change how symmetric a face is. So if the detector gives different
geometry for the photo and its mirror image, the detector is unsure, and we can measure
that without any ground truth.''
\end{say}

\begin{careful}
Flip test-time augmentation already exists in human-pose estimation, where it is used to
improve accuracy. Do not claim flipping itself is new. What is new here is using it, for a
\emph{symmetry} measurement, as a reliability signal and to cancel left/right bias.
\end{careful}

\subsection{Step 6: the quality and pose gate (\texttt{quality.py})}
The gate looks at seven things the system can actually measure:
\begin{center}\small
\begin{tabular}{@{}p{3.4cm}p{6.6cm}p{5.4cm}@{}}
\toprule
Signal & How it is computed & Why it matters\\
\midrule
Face size & IOD in pixels & tiny faces: 1 pixel is a large \% of IOD\\
Sharpness & $\log_{10}$ variance of the Laplacian of the face crop at 256\,px width & blur moves landmarks\\
Brightness (low/high) & mean grey level of the face crop & dark or burnt-out faces\\
Clipped pixels & share of crop pixels near black or white & lost detail\\
$|$yaw$|$ & from the transformation matrix & foreshortens one side\\
$|$pitch$|$ & from the transformation matrix & distorts proportions\\
Mirror inconsistency $\delta$ & Step 5 & detector unsure\\
\bottomrule
\end{tabular}
\end{center}
\textbf{Decision rule:} each signal has a caution limit and a reject limit; the limits are
stored in the file \texttt{configs/gate\_thresholds.json}. The \emph{worst} signal decides the
status, and every triggered signal is listed with a plain-language reason. A rejected photo
shows no score unless the user explicitly asks to see it, and it is then labelled unreliable.

\textbf{How the limits were set (important for the viva):}
\begin{itemize}
\item \emph{A priori}, from the measurement's assumptions: yaw 10$^\circ$/20$^\circ$, pitch
15$^\circ$/25$^\circ$, IOD 60/30\,px.
\item \emph{From the training split}: percentiles of the 810 training faces (2.5th
percentile for sharpness and brightness; 90th/99th for $\delta$; 97.5th for region validity).
\item The test split was \textbf{never} used to choose any limit.
\end{itemize}
\textbf{Roll is not gated} because the fitted axis already handles it.

\subsection{Step 7: the readable index}
The app also shows $S = 100/(1+\alpha A)$ with $\alpha$ = \alphaIndex{}, chosen so that the
median training face (A = \trainMedian\% IOD) scores 50. It is only a friendlier way to show
$A$ (higher = more symmetric). \textbf{It is not a validated clinical or beauty scale.}

\subsection{The four variants we compare}
\begin{center}
\begin{tabular}{@{}lp{11.5cm}@{}}
\toprule
M0 & Baseline: KD-tree matching, vertical midline, align + re-detect, / face width\\
M1 & Fixed pairs + 68-point map + fitted axis, / IOD (no gate)\\
M2 & M1 + quality/pose gate (without the mirror test)\\
M3 & Full system: mirror-averaged landmarks, gate including $\delta$, region validity\\
\bottomrule
\end{tabular}
\end{center}
All four are computed from the \emph{same detections} of each image, so any difference comes
from the method, not from detector randomness.

%======================================================================
\section{The data}
%======================================================================
\textbf{Dataset: LFPW} (Labelled Face Parts in the Wild), part of the 300-W benchmark. Each
image has 68 landmarks placed by human experts. We used the copy published on GitHub by the
Menpo project (\texttt{github.com/menpo/lfpw-train}, \texttt{lfpw-test}).
\begin{itemize}
\item \textbf{Train split} (811 images, 810 faces detected): used \emph{only} to fit the 478$\to$68
map and to set gate limits.
\item \textbf{Test split} (224 images, all detected): every reported number.
\item Each photo is cropped around the annotated face first, as in the 300-W protocol.
Many LFPW photos contain several people, and without the crop the detector sometimes
analysed the wrong person (this initially caused a 32\% mean landmark error, which dropped
to \nmeFitInner\% after the crop).
\end{itemize}
\textbf{Why not WFLW?} It was the first choice (it has pose, blur and occlusion subsets), but
its images are hosted on Google Drive, which our compute environment could not reach. The
code reads any 68-point \texttt{.pts} dataset, so WFLW or other 300-W subsets can be added later.\\
\textbf{Licences and privacy:} 300-W is for research use only. No dataset image is stored
in our repository or shown in the paper or slides; figures draw landmark geometry only.
